% ========== EXTENSION LIFECYCLE "CHAMBERING" ==========

\section{Extension Lifecycle: "Chambering"}
\label{sec:extension-lifecycle-chambering}

\lettrine{T}{he} extension lifecycle in Symphony is managed through a sophisticated system called \textbf{"Chambering"}, which orchestrates the loading, execution, and unloading of extensions with precise state management and resource control. The term "Chambering" reflects the musical metaphor of preparing instruments for performance—extensions are carefully prepared, positioned, and coordinated before they contribute to the development symphony.

\subsection{Lifecycle States}
\label{subsec:lifecycle-states}

The Chambering system defines seven distinct lifecycle states that extensions progress through during their operational lifetime. Each state represents a specific phase of extension readiness and capability.

\begin{figure}[htbp]
\centering
\begin{tikzpicture}[
    node distance=2cm,
    state/.style={rectangle, rounded corners, draw=brandPrimary, fill=brandPrimary!10, text width=2.5cm, text centered, minimum height=1cm},
    transition/.style={->, >=stealth, thick, brandSecondary}
]

% States
\node[state] (dormant) {Dormant};
\node[state, right of=dormant] (loading) {Loading};
\node[state, right of=loading] (initializing) {Initializing};
\node[state, below of=initializing] (ready) {Ready};
\node[state, left of=ready] (active) {Active};
\node[state, left of=active] (suspended) {Suspended};
\node[state, below of=active] (unloading) {Unloading};

% Transitions
\draw[transition] (dormant) -- (loading) node[midway, above] {\small Load};
\draw[transition] (loading) -- (initializing) node[midway, above] {\small Init};
\draw[transition] (initializing) -- (ready) node[midway, right] {\small Complete};
\draw[transition] (ready) -- (active) node[midway, above] {\small Activate};
\draw[transition] (active) -- (suspended) node[midway, above] {\small Suspend};
\draw[transition] (suspended) -- (active) node[midway, below] {\small Resume};
\draw[transition] (active) -- (unloading) node[midway, right] {\small Unload};
\draw[transition] (ready) to[bend left=45] (unloading) node[midway, below] {\small Cleanup};

% Error transitions (dashed)
\draw[transition, dashed, red] (loading) to[bend right=30] (unloading) node[midway, below left] {\small Error};
\draw[transition, dashed, red] (initializing) to[bend right=15] (unloading) node[midway, below] {\small Failure};

\end{tikzpicture}
\caption{Extension Lifecycle State Diagram}
\label{fig:extension-lifecycle-states}
\end{figure}

\subsubsection{State Definitions}
\label{subsubsec:state-definitions}

\begin{description}[leftmargin=3cm,labelwidth=2.5cm]
\item[\textbf{Dormant}] Extension exists in the registry but is not loaded into memory. Minimal resource footprint.

\item[\textbf{Loading}] Extension binary/code is being loaded into the appropriate execution context (Pit or Grand Stage).

\item[\textbf{Initializing}] Extension is setting up its runtime environment, establishing connections, and preparing for operation.

\item[\textbf{Ready}] Extension is fully initialized and capable of handling requests but not actively processing.

\item[\textbf{Active}] Extension is actively processing requests and contributing to development workflows.

\item[\textbf{Suspended}] Extension is temporarily paused, maintaining state but not processing new requests.

\item[\textbf{Unloading}] Extension is cleaning up resources and preparing for removal from memory.
\end{description}

\subsection{State Transitions \& Validation}
\label{subsec:state-transitions-validation}

State transitions in the Chambering system are governed by strict validation rules that ensure system stability and resource integrity. Each transition includes pre-conditions, validation checks, and post-conditions.

\begin{table}[h]
\centering
\caption{State Transition Validation Matrix}
\label{tab:state-transition-validation}
\begin{tabular}{@{}llll@{}}
\toprule
\textbf{Transition} & \textbf{Pre-conditions} & \textbf{Validation} & \textbf{Timeout} \\
\midrule
Dormant → Loading & Manifest valid & Dependency check & 30s \\
Loading → Initializing & Binary loaded & Security scan & 60s \\
Initializing → Ready & Init complete & Health check & 45s \\
Ready → Active & Resources available & Capability test & 15s \\
Active → Suspended & Safe state & State preservation & 10s \\
Suspended → Active & Resources available & State restoration & 15s \\
* → Unloading & Any state & Cleanup validation & 30s \\
\bottomrule
\end{tabular}
\end{table}

\subsubsection{Transition Triggers}
\label{subsubsec:transition-triggers}

State transitions are triggered by various events and conditions within the Symphony ecosystem:

\begin{expandedlist}
\item \textbf{User Actions}: Explicit extension activation/deactivation through UI
\item \textbf{Conductor Decisions}: AI-driven orchestration based on context and performance
\item \textbf{System Events}: File operations, project changes, or workflow triggers
\item \textbf{Resource Constraints}: Memory pressure, CPU limits, or performance thresholds
\item \textbf{Dependency Changes}: Updates to required extensions or system components
\item \textbf{Error Conditions}: Failures, timeouts, or security violations
\end{expandedlist}

\subsection{Lifecycle Hooks}
\label{subsec:lifecycle-hooks}

The Chambering system provides lifecycle hooks that allow extensions to participate in state management and coordinate with the broader Symphony ecosystem.

\subsubsection{Standard Lifecycle Hooks}
\label{subsubsec:standard-lifecycle-hooks}

\begin{table}[h]
\centering
\caption{Extension Lifecycle Hooks}
\label{tab:lifecycle-hooks}
\begin{tabular}{@{}lll@{}}
\toprule
\textbf{Hook} & \textbf{Timing} & \textbf{Purpose} \\
\midrule
\texttt{on\_load} & Before Loading → Initializing & Resource preparation \\
\texttt{on\_init} & During Initializing & Setup and configuration \\
\texttt{on\_ready} & Initializing → Ready & Final preparations \\
\texttt{on\_activate} & Ready → Active & Begin processing \\
\texttt{on\_suspend} & Active → Suspended & Pause operations \\
\texttt{on\_resume} & Suspended → Active & Resume operations \\
\texttt{on\_deactivate} & Active → Ready & Stop processing \\
\texttt{on\_unload} & Before Unloading & Cleanup and shutdown \\
\bottomrule
\end{tabular}
\end{table}

\subsubsection{Hook Implementation Example}
\label{subsubsec:hook-implementation-example}

Extensions implement lifecycle hooks through their manifest and runtime interface:

\begin{lstlisting}[language=toml,caption=Lifecycle Hook Configuration in Symphony.toml]
[lifecycle]
hooks = [
    "on_load",
    "on_init", 
    "on_activate",
    "on_suspend",
    "on_unload"
]

[lifecycle.timeouts]
on_init = 30000      # 30 second initialization timeout
on_activate = 5000   # 5 second activation timeout
on_suspend = 10000   # 10 second suspension timeout
on_unload = 15000    # 15 second cleanup timeout

[lifecycle.error_handling]
retry_on_failure = true
max_retries = 3
backoff_strategy = "exponential"
\end{lstlisting}

\subsection{Resource Management \& Cleanup}
\label{subsec:resource-management-cleanup}

The Chambering system implements sophisticated resource management to ensure system stability and optimal performance across the extension ecosystem.

\subsubsection{Resource Tracking}
\label{subsubsec:resource-tracking}

Symphony tracks multiple resource categories for each extension:

\begin{compactlist}
\item \textbf{Memory}: Heap allocation, stack usage, shared memory regions
\item \textbf{CPU}: Processing time, thread count, scheduling priority
\item \textbf{I/O}: File handles, network connections, disk usage
\item \textbf{System}: Process handles, registry entries, temporary files
\item \textbf{Extension-Specific}: AI model memory, cache storage, persistent data
\end{compactlist}

\begin{infobox}[title=Resource Quotas by Extension Type]
\textbf{Instruments}: High memory (512MB-2GB), moderate CPU, network access \\
\textbf{Operators}: Moderate memory (64MB-256MB), high CPU, limited I/O \\
\textbf{Motifs}: Low memory (32MB-128MB), low CPU, UI resources only
\end{infobox}

\subsubsection{Cleanup Strategies}
\label{subsubsec:cleanup-strategies}

The system employs multiple cleanup strategies based on extension type and resource usage patterns:

\begin{figure}[htbp]
\centering
\begin{tikzpicture}[
    box/.style={rectangle, draw=brandPrimary, fill=brandPrimary!10, text width=3cm, text centered, minimum height=1.5cm},
    arrow/.style={->, >=stealth, thick, brandSecondary}
]

% Cleanup phases
\node[box] (immediate) {Immediate Cleanup};
\node[box, right=2cm of immediate] (deferred) {Deferred Cleanup};
\node[box, right=2cm of deferred] (background) {Background Cleanup};

% Arrows
\draw[arrow] (immediate) -- (deferred) node[midway, above] {\small 5s delay};
\draw[arrow] (deferred) -- (background) node[midway, above] {\small 30s delay};

% Details below each phase
\node[below=0.5cm of immediate, text width=3cm, text centered] {
    \small
    • Close handles \\
    • Release locks \\
    • Stop threads
};

\node[below=0.5cm of deferred, text width=3cm, text centered] {
    \small
    • Free memory \\
    • Clear caches \\
    • Save state
};

\node[below=0.5cm of background, text width=3cm, text centered] {
    \small
    • Delete temp files \\
    • Compact logs \\
    • Update metrics
};

\end{tikzpicture}
\caption{Three-Phase Resource Cleanup Strategy}
\label{fig:cleanup-strategy}
\end{figure}

\subsubsection{Memory Management}
\label{subsubsec:memory-management}

Symphony implements advanced memory management techniques tailored to different extension types:

\begin{table}[h]
\centering
\caption{Memory Management by Extension Type}
\label{tab:memory-management}
\begin{tabular}{@{}llll@{}}
\toprule
\textbf{Extension Type} & \textbf{Allocation Strategy} & \textbf{Cleanup Policy} & \textbf{Monitoring} \\
\midrule
Instruments & Pool-based allocation & Lazy cleanup & Continuous \\
Operators & Stack-preferred & Immediate cleanup & Periodic \\
Motifs & Shared UI pools & Reference counting & Event-driven \\
\bottomrule
\end{tabular}
\end{table}

\subsection{Performance Optimization}
\label{subsec:performance-optimization}

The Chambering system includes several performance optimization strategies that adapt to usage patterns and system conditions.

\subsubsection{Predictive Loading}
\label{subsubsec:predictive-loading}

The Conductor analyzes usage patterns to predict which extensions will be needed and pre-loads them during idle periods:

\begin{successbox}
\textbf{Smart Pre-loading:} Extensions frequently used together are loaded as groups, reducing activation latency by up to 80\% for common workflows.
\end{successbox}

\subsubsection{Adaptive Resource Allocation}
\label{subsubsec:adaptive-resource-allocation}

Resource limits adjust dynamically based on system load and extension behavior:

\begin{compactlist}
\item High-performing extensions receive increased resource quotas
\item Problematic extensions face stricter limits and monitoring
\item System-wide resource pressure triggers automatic suspension of low-priority extensions
\item Performance metrics influence future loading decisions
\end{compactlist}

The Chambering system represents a sophisticated approach to extension lifecycle management that balances performance, security, and resource efficiency while providing the flexibility needed for Symphony's AI-first development environment.